\chapter{Introduction}
\label{ch:introduction}

\section{Aim}
\label{sec:introduction:aim}

This project provides a framework for developers building exchange systems.
Applications built on the framework are event-driven. Their reactor dispatches
four kinds of events to application threads: topic-based publish and subscribe
records, inter-thread communication (ITC) messages, timer events, and activity
on unicast socket connections. Topic-based publish and subscribe is a central
framework facility. The framework is designed for low latency, not
ultra-low latency.

The framework also provides high-availability facilities for exchange
applications. Paired components use leader and follower roles, and leases
decide which instance may act. This is a simpler design than one based on a
general consensus algorithm such as Raft or Paxos
\cite{ongaro2014raft,lamport2001paxos}. The book states the
behaviour required of the framework and the example venue, including the
conditions and assumptions on which that behaviour depends.

\section{Topic-based publish and subscribe}
\label{sec:introduction:pubsub}

The framework provides topic-based publish and subscribe for applications
that need to distribute ordered records to independent subscribers. A
subscriber can resume from an earlier position to catch up before receiving
new records. The publisher retains the history in its write-ahead log and
sends each subscriber a copy over that subscriber's own unicast TCP
connection.

This delivery design reflects the lack of a reliable UDP multicast system
that has gained broad industry adoption. Protocols such as PGM and NORM have
been specified for decades \cite{rfc3208,rfc5740}, but have not become widely
adopted solutions. Other designs, including redundant A and B channels, were
considered and rejected. Unicast fanout makes delivery to each subscriber
independent, at the cost of work and network traffic that grow with the number
of subscribers. The framework's delivery and recovery requirements are set
out in \cref{sec:pubsub}.

\section{Latency and how it is measured}
\label{sec:introduction:latency}

The design aims to keep work on the order path bounded and avoid delays caused
by waiting for other threads. The framework uses event-driven dispatch,
lock-free queues for inter-thread communication, single-thread ownership of
changing state, and pre-allocated memory on hot paths. These choices reduce
work and avoid common sources of variable delay; they do not make latency
constant or guarantee a particular response time.

The venue measures an order's round trip from when its gateway reads the
order from the member's connection until the first byte of the execution
report is sent back. It also records measurements at stages of the order path
and on reactor and inter-thread hand-offs. Controlled load runs hold the
offered rate and the load client's processor placement steady, then compare
the resulting distributions across runs. The number of observations matters:
the median and common percentiles can be compared with more confidence than
rare tail percentiles from a small run. These are measurements of the tested
configuration on a particular machine, not guarantees for every deployment.
The measurement method and current findings are in
\texttt{docs/operations/latency\_findings.md}.

\section{The example applications}
\label{sec:introduction:applications}

The project includes example exchange applications to show how the framework
can be used. They include FIX and binary order gateways, a sequencer, a
matching engine, a matching engine publisher, and services for authentication,
administration, arbitration, and witnessing. They are examples, not a
fully-fledged exchange. For instance, the matching engine accepts and cancels
orders but does not match buyers and sellers. The examples exercise framework
facilities under exchange-like traffic and failure conditions; their
limitations are stated where relevant in the book.

Two Java web applications support the examples. The FIX test client lets a
developer or operator connect through either the FIX or binary gateway, send
and cancel orders, inspect reports, and drive repeatable tests with scripts.
The administration service manages firms, member comp IDs, gateway
permissions, and credentials. They are supporting applications: the test
client exercises the venue, while the administration service provisions
access to it.

\section{How to read this book}
\label{sec:introduction:structure}

This is a functional specification, not a source-code or configuration
reference. Its requirements state observable behaviour the framework or venue
must provide, explain why each requirement matters, and identify the tests
that verify it. Where the implementation does not meet a requirement, the
book records the gap rather than changing the requirement to describe the
shortfall. Deliberate limitations are also stated so that an example's
boundaries are clear.

The chapters follow the system from its foundations to its failure and
recovery behaviour:

\begin{itemize}[nosep]
  \item \textbf{The framework} states the requirements for the reactor,
        threads, inter-thread and inter-process communication, memory,
        serialisation, timers, the write-ahead log, configuration, logging,
        and topic-based publish and subscribe.
  \item \textbf{The applications built on the library} states the ordinary
        behaviour of the example exchange components and the services around
        them.
  \item \textbf{Building, configuring, releasing and deploying} describes
        the project workflow, TOML configuration, and the environment
        profiles used for development, testing, pre-production, and production.
  \item \textbf{Security} states the requirements for protecting identities,
        credentials, and communications.
  \item \textbf{High availability} states what the venue must preserve when
        processes or machines fail, and how components recover.
  \item \textbf{Metrics} states what the system measures and makes available
        for operational observation.
  \item \textbf{Disaster recovery} addresses site loss and site switchover:
        the decisions and procedures needed to resume trading elsewhere.
\end{itemize}

Implementation details, message layouts, configuration keys, and operating
instructions are documented separately and linked from the relevant
requirements where they help explain the system.
